% Generated draft: requires mathematical and source review.
\documentclass[11pt]{article}
\usepackage[margin=25mm]{geometry}
\usepackage{fontspec}
\setmainfont{Latin Modern Roman}
\usepackage{amsmath,amssymb,mathtools,unicode-math}
\setmathfont{Latin Modern Math}
\usepackage{enumitem,needspace,xurl,hyperref}
\setlength{\parindent}{0pt}
\setlength{\parskip}{5pt}
\setlength{\emergencystretch}{3em}
\newcommand{\N}{\mathbb{N}}
\newcommand{\Z}{\mathbb{Z}}
\newcommand{\Q}{\mathbb{Q}}
\newcommand{\R}{\mathbb{R}}
\newcommand{\C}{\mathbb{C}}
\newcommand{\F}{\mathbb{F}}
\newcommand{\A}{\mathbb{A}}
\newcommand{\PP}{\mathbb{P}}
\newcommand{\E}{\mathbb{E}}
\newcommand{\eps}{\varepsilon}
\newcommand{\dd}{\,\mathrm{d}}
\begin{document}
\section*{UnsolvedMath ID: 30002695}
\textbf{Source problem code:} catalogue entry
\subsection{Definitions and mathematical statement}
Let \((M,g)\) be a homogeneous Riemannian manifold, meaning that its isometry group acts transitively on \(M\). Write \(R_g\) for its Riemann curvature tensor and \(\nabla\) for the Levi--Civita connection. With the usual formal-adjoint convention, the rough Laplacian on tensors is
\[
\nabla^*\nabla R_g=-\operatorname{tr}_g(\nabla^2R_g).
\]
The supplied record names \(k_o\), \(k_t\), and \(k_S\) as the curvature-derivative, traceless-curvature, and Singer invariants, respectively. It supplies neither their definitions nor the defining traceless-curvature relation referred to as Equation (5). These names alone do not determine the conventions needed for a complete mathematical formulation.

Subject to recovering those definitions and Equation (5), the supplied questions are:
\begin{enumerate}
\item For every homogeneous Riemannian manifold for which \(k_o\) is finite, must \(k_o\) be an even integer?
\item Under the hypothesis \(k_t=1\), what conclusions hold concerning the equation
\[
\nabla^*\nabla R_g=\lambda R_g?
\]
The record does not specify whether this asks for existence of a constant \(\lambda\), restrictions on its possible values, or a classification of metrics satisfying the equation. The description as an eigenvalue equation suggests a constant spectral parameter, but that interpretation is not established by the supplied definitions.
\item For a Ricci flow \(g(t)\) starting from a homogeneous metric, how does Equation (5) behave, and how do
\[
k_t(g(t)),\qquad k_o(g(t)),\qquad k_S(g(t))
\]
vary with time? In particular, no specific evolution identity or preservation assertion is supplied as the intended target. The normalization of the flow and its time interval are unspecified. The usual unnormalized convention is \(\partial_tg=-2\operatorname{Ric}_g\), but the supplied record does not establish that this is the intended convention.
\end{enumerate}
The parity question occurs in the supplied clean statement but not in the supplied original statement. Equation (5) must be recovered before its evolution can be formulated explicitly. Thus this is a qualified statement of the recorded questions, requiring human mathematical review.
\subsection{Short English statement}
For homogeneous Riemannian manifolds, the record asks whether finite \(k_o\) is always even, what \(k_t=1\) implies about a curvature rough-Laplacian eigenvalue equation, and how three curvature invariants and an omitted defining equation behave under Ricci flow. Missing definitions and Equation (5) prevent a fully precise formulation.
\subsection{Sources}
\begin{itemize}
\item Supplied local UnsolvedMath registry record: permanent ID \(30002695\), source problem code \texttt{catalogue entry}, title \emph{Parity and Ricci-Flow Evolution of Curvature Invariants}; collection: Oberwolfach Reports Open Problems. No individual proposer is supplied; the proposed year is recorded as 2014.
\item Supplied report citation: \emph{Mini-Workshop: Einstein Metrics, Ricci Solitons and Ricci Flow under Symmetry Assumptions} (2015), Oberwolfach Reports, pp.\ 2529--2568, DOI \texttt{10.4171/owr/2014/45}. Supplied source URL and literature-source link, labeled ``Oberwolfach Report 45/2014'': \url{https://doi.org/10.4171/owr/2014/45}. The citation gives 2015 while the report label and proposed year give 2014; this discrepancy is unresolved here.
\item Literature assessment dated exactly \texttt{2026-08-21}: the supplied curation labels the problem open and reports that no general resolution was identified for questions described there as concerning an additional inequality and the Ricci-flow evolution of geometric invariants. This is supplied dated curation, not an independently verified literature assessment. The background explicitly requests independent verification of time-sensitive status and cited claims.
\item The registry labels the statement \texttt{corrected\_verified} and reports comparison with the official Oberwolfach report. These are supplied verification claims; the report has not been independently inspected for this draft.
\end{itemize}
\end{document}
